% !TeX program = xelatex
% 本模板由 Bryce 制作，用于浙江大学海宁校区军事理论课程。
% 本文件包含填写提示和排版示例；提交前请替换全部示例内容。
\documentclass{milpaper}
\addbibresource{references.bib} % 文献资料在此文件填写，格式由国标样式自动生成。

% ===== 1. 基本信息：通常只需编辑 main.tex，无需修改 milpaper.cls =====
\newcommand{\authorname}{姓名}
\newcommand{\studentnumber}{学号}
\newcommand{\authorenglish}{YOUR NAME}
\newcommand{\papertitle}{论文中文题目}
\newcommand{\papersubtitle}{论文副标题} % 无副标题时改为 {}
\newcommand{\papertitleenglish}{English Title of the Paper}

\begin{document}
\paperTitle{\papertitle}{\papersubtitle}

% ===== 2. 中英文摘要与关键词 =====
\begin{cnabstract}
【请填写中文摘要。简要介绍研究问题、使用的材料或分析方法、主要论点及结论。摘要应独立成段，不要直接照抄引言，也不要保留本段填写提示。摘要长度以任课教师要求为准。】
\end{cnabstract}
\cnkeywords{关键词一；关键词二；关键词三}

\englishTitle{\papertitleenglish}
\begin{enabstract}
[Replace this paragraph with your English abstract. Briefly state the research question, materials or methods, main argument, and conclusion. Check that it conveys the same meaning as the Chinese abstract.]
\end{enabstract}
\enkeywords{Keyword One; Keyword Two; Keyword Three}

% ===== 3. 正文：以下为结构及排版示例，可按选题自由增删 =====
\section{引言}
【在此说明选题背景、拟讨论的问题及其意义。结合课程内容限定讨论范围，明确论文准备回答什么问题。】

【介绍材料来源与文章结构。引用他人观点、数据或图表时应标注出处\cite{example-journal-cn,example-journal-en}。】

\section{概念与材料}
\subsection{核心概念}
【解释核心概念，说明其与研究问题的关系。段首自动缩进。】

\subsubsection{分析范围}
【限定研究对象、时间范围和比较维度，避免标题层级过多。】

\subsection{材料选择与核验}
【介绍资料类型、来源及局限，核验事实并注明出处\cite{example-book,example-web}。】

\section{分析与讨论}
\subsection{论证结构与插图示例}
【围绕论点展开分析，区分事实与推论。图\ref{fig:workflow}仅演示插图及交叉引用。】

% 下图使用 TikZ 绘制，无需额外图片文件。
\begin{figure}[htbp]
  \centering
  \begin{tikzpicture}[
    paperstage/.style={draw=black!65,rectangle,minimum width=29mm,
      minimum height=10mm,align=center,font=\small},
    flow/.style={-{Stealth[length=2mm]},line width=.55pt}]
    \node[paperstage] (question) {提出问题};
    \node[paperstage,right=7mm of question] (sources) {查阅资料};
    \node[paperstage,right=7mm of sources] (argument) {形成论证};
    \node[paperstage,right=7mm of argument] (revision) {检查修订};
    \draw[flow] (question)--(sources);
    \draw[flow] (sources)--(argument);
    \draw[flow] (argument)--(revision);
  \end{tikzpicture}
  \caption{论文写作流程示意（排版示例）}
  \label{fig:workflow}
  \sourceNote{资料来源：模板绘制，仅演示排版。}
\end{figure}

% 使用自己的图片时，可将上面的 tikzpicture 环境替换为：
% \includegraphics[width=.85\textwidth]{images/figure.png}
% 请先将图片上传到项目的 images 文件夹，并修改图题、标签与来源。

\subsection{比较分析与表格示例}
【按研究问题选择比较维度，并注明来源。表\ref{tab:example}仅为三线表排版示例。】

\begin{table}[htbp]
  \centering
  \caption{比较分析表（排版示例）}
  \label{tab:example}
  \zihao{-5}
  \renewcommand{\arraystretch}{1.2}
  \begin{tabularx}{\textwidth}{@{}>{\heiti\raggedright\arraybackslash}p{25mm}
    >{\raggedright\arraybackslash}X>{\raggedright\arraybackslash}X@{}}
    \toprule
    比较维度 & 对象一 & 对象二 \\
    \midrule
    背景与条件 & 【填写相应材料】 & 【填写相应材料】 \\
    主要特点 & 【填写分析结果】 & 【填写分析结果】 \\
    局限与启示 & 【填写论证与边界】 & 【填写论证与边界】 \\
    \bottomrule
  \end{tabularx}
  \sourceNote{资料来源：【填写实际来源；如自行整理，请注明所依据的文献。】}
\end{table}

\section{结论}
【回答研究问题，概括主要发现与依据，说明材料或方法的局限。】

% ===== 4. 参考文献：全部替换为实际阅读和引用的文献 =====
% 在 references.bib 中填写文献；正文用 \cite{引用键} 引用。
% 采用 GB/T 7714-2025 顺序编码制；只列出正文实际引用的条目。
% 超过三位作者时自动列前三位，中文加“等”，英文加斜体 et al.
% 示例文献均为虚构格式占位，不能直接提交。
\printbibliography

% ===== 5. 课程心得与附录：按任课教师要求保留、调整或删除 =====
\subsection*{课程（论文）心得与建议:}
【填写本人真实的课程学习或论文写作心得，以及对课程的意见和建议。可选。】

\Needspace{10\baselineskip}
\section*{附录\quad GAI使用说明与相关截图}
【按课程要求如实说明 GAI 工具、用途、参与环节及核验过程。可选】

% 如需截图，将下列占位框替换为真实图片，例如：
% \begin{center}
%   \includegraphics[width=.9\textwidth,height=.55\textheight,
%     keepaspectratio]{images/gai-screenshot.png}
% \end{center}
% 分享或提交截图前，请检查其中的姓名、账号、头像与其他无关个人信息。
\begin{center}
  \fbox{\parbox[c][20mm][c]{.88\textwidth}{\centering
    【按课程要求插入实际使用过程的截图】\\[4pt]
    无需截图时，请删除此占位框。
  }}
\end{center}

\end{document}
